\section{The scalar square: theorem and physics}
\label{sec:scalar-square-status}

Let
$$X=S\times E,$$
where $S$ is a nonsingular projective K3 surface and $E$ is an elliptic
curve.  Fix primitive classes
$$\beta_h\in H_2(S,\Z),\qquad \langle \beta_h,\beta_h\rangle=2h-2,$$
and write
$$(\beta_h,d)=\iota_{S*}\beta_h+\iota_{E*}(d[E])\in H_2(X,\Z).$$
All reciprocal Siegel-form identities below are identities of Laurent
expansions in the Fourier chamber fixed by the Borcherds product.
Constants $C_*$ are nonzero complex scalars; they absorb only the
normalization of \(\chi_{10}^{\mathrm{OP}}\) versus \(\Delta_5^2\) and
never a nonconstant factor.

Put
$$\Delta_{10}^{\theta}(Z):=\Delta_5(Z)^2,\qquad
\chi_{10}^{\mathrm{OP}}:=2^{-12}\Delta_{10}^{\theta}.$$

\begin{definition}[Oberdieck--Pixton primitive reduced series]
Let $\beta_h^\vee\in H^2(S,\Q)$ satisfy
$$\langle \beta_h,\beta_h^\vee\rangle=1.$$
Define the disconnected reduced Gromov--Witten invariant
$$
\mathsf N^X_{g,h,d}
=
\int_{[\overline M^{\bullet}_{g,1}(X,(\beta_h,d))]^{\mathrm{red}}}
\operatorname{ev}_1^*
\left(\pi_1^*\beta_h^\vee\cup \pi_2^*[\mathrm{pt}]\right).
$$
The primitive reduced curve-counting partition function is
$$
Z^X_{\mathrm{OP}}(u,Q,T)
=\sum_{g\geq0}\sum_{h\geq0}\sum_{d\geq0}
\mathsf N^X_{g,h,d}\,u^{2g-2}Q^{h-1}T^{d-1}.
$$
With $P=e^{iu}$, the notation $Z^X_{\mathrm{OP}}(P,Q,T)$ means this
Laurent expansion under the change of variables $P=e^{iu}$.
\end{definition}

\begin{theorem}[Oberdieck--Pixton primitive theorem]
\label{thm:op-primitive-square-status}
Oberdieck--Pixton prove the primitive Igusa cusp-form conjecture
\cite{OB}; this is Theorem~1 of their paper:
$$
Z^X_{\mathrm{OP}}(P,Q,T)
=-\frac{1}{\chi_{10}^{\mathrm{OP}}(P,T,Q)}.
$$
Equivalently, after the fixed normalization
$$\chi_{10}^{\mathrm{OP}}(P,T,Q)=2^{-12}\Delta_5(Z(P,Q,T))^2,$$
one has the scalar square
$$
Z^X_{\mathrm{OP}}(P,Q,T)
=\frac{C_{\mathrm{OP}}}{\Delta_5(Z(P,Q,T))^2},
\qquad C_{\mathrm{OP}}\in\C^\times.
$$
This is the primitive reduced curve-counting theorem.  The
Hilbert-scheme DT equality below uses it after the reduced DT/PT
comparison and the primitive GW/PT comparison have fixed the
quotient-by-\(E\) normalization.
\end{theorem}

\begin{definition}[Reduced Hilbert-scheme DT series]
Let
$$
\operatorname{Hilb}^n(X,(\beta_h,d))
=
\{Y\subset X\mid [Y]=(\beta_h,d),\ \chi(\mathcal O_Y)=n\}.
$$
The translation action of $E$ on the second factor acts on this Hilbert
scheme.  Define the Behrend-weighted quotient invariant
$$
\mathbf{DT}^{X,\mathrm{Hilb}}_{n,h,d}
=
\int_{\operatorname{Hilb}^n(X,(\beta_h,d))/E}\nu_{\mathrm{Hilb}}\,de,
$$
where $\nu_{\mathrm{Hilb}}$ is Behrend's constructible function
\cite{BEHREND} and $de$ denotes orbifold topological Euler
characteristic.  The corresponding Hilbert-scheme DT partition function
in the manuscript normalization is
$$
Z^X_{\mathrm{DT}}(T,Q,P)
=
\sum_{h\geq0}\sum_{d\geq0}\sum_{n\in\Z}
\mathbf{DT}^{X,\mathrm{Hilb}}_{n,h,d}\,
T^{d-1}Q^{h-1}(-P)^n.
$$
This definition is the Donaldson--Thomas normalization attached to
ideal sheaves on $X$ \cite{DT}.  The equality with \(\Delta_5^{-2}\)
requires the comparison theorem below.
\end{definition}

\begin{theorem}[Reduced Hilbert-scheme DT scalar square]
\label{thm:hilb-dt-scalar-normalization-status}
In the quotient-by-\(E\) normalization,
$$
Z^X_{\mathrm{DT}}(T,Q,P)
=Z^X_{\mathrm{OP}}(P,Q,T)
=-4096\,\Delta_5(Z)^{-2}.
$$
\end{theorem}

\begin{proof}
Oberdieck's reduced DT/PT comparison for \(S\times E\),
\cite{OberdieckReducedPT}*{Theorem 3(i)}, gives equality of reduced
Hilbert-scheme DT and stable-pairs series for every non-zero curve
class.  The degree-zero factor is \(M(-P)^{\chi(S\times E)}=1\).  For
primitive \(\beta_h\), the GW/PT comparison used in
\cite{OB}*{Section 7.1, equation (GWPTcorr)} and
\cite{OPand}*{Proposition 5} identifies the stable-pairs series with
the Oberdieck--Pixton reduced series after \(P=e^{iu}\).  The
Oberdieck--Pixton theorem then gives the displayed Igusa reciprocal.
\end{proof}

\begin{remark}[Low-degree checks]
Bryan's topological-vertex computation of the Behrend-weighted
Hilbert-scheme theory for primitive classes of square $-2$ and $0$
\cite{BRYAN} gives low-degree checks of the same quotient-by-\(E\)
normalization.
\end{remark}

\begin{definition}[Scalar-square branch]
\label{def:scalar-square-branch-status}
A scalar-square branch is the datum of a source
$$\mathfrak s\in\{\mathrm{OP},\mathrm{DT}\}$$
and a partition function $Z^X_{\mathfrak s}$ satisfying
$$
Z^X_{\mathfrak s}(Z)=\frac{C_{\mathfrak s}}{\Delta_5(Z)^2},
\qquad C_{\mathfrak s}\in\C^\times.
$$
The unconditional branches are
$$
Z^X_{\square,\mathrm{uncond}}:=Z^X_{\mathrm{OP}},
\qquad
Z^X_{\square,\mathrm{DT}}:=Z^X_{\mathrm{DT}},
$$
$$
C_{\square,\mathrm{uncond}}=C_{\square,\mathrm{DT}}=-4096.
$$
If the source is not displayed, $Z^X_\square$ denotes an arbitrary
branch satisfying the displayed scalar-square identity; it must not be
identified with a protected D-brane trace without the conjectural
physics dictionary below.
\end{definition}

\begin{theorem}[Square-root consequence]
\label{thm:scalar-square-consequence-status}
For every scalar-square branch
$$Z^X_\square=\frac{C_\square}{\Delta_5^2},$$
and for the Borcherds determinant
$$\mathcal D_X(Z):=\Delta_5(Z),$$
one has
$$
\mathcal D_X(Z)^2=\frac{C_\square}{Z^X_\square(Z)}.
$$
With
$$\mathcal Z_{\mathrm{BPS}}^{\mathrm{vir}}(Z):=\mathcal D_X(Z)^{-1},$$
this is equivalently
$$
Z^X_\square(Z)
=C_\square\left(\mathcal Z_{\mathrm{BPS}}^{\mathrm{vir}}(Z)\right)^2.
$$
\end{theorem}

\begin{proof}
The branch identity gives
$$C_\square/Z^X_\square=\Delta_5^2.$$
The chosen square root is not arbitrary: it is the Borcherds product
already denoted by $\mathcal D_X=\Delta_5$.  Inverting gives
$$Z^X_\square=C_\square(\mathcal D_X^{-1})^2.$$
\end{proof}

\begin{conjecture}[Protected-index dictionary]
\label{conj:protected-index-dictionary-status}
For D6--D2--D0 charge
$$\gamma(h,d,n)=(1;\,(\beta_h,d);\,n),$$
the physical protected index is expected to be represented by the
Behrend-weighted Hilbert-scheme DT invariant:
$$
\Omega_X(h,d,n)
=\operatorname{Tr}_{\mathcal H_{\gamma(h,d,n)}}(-1)^F
\ \stackrel{\mathrm{phys}}{=}\ 
\mathbf{DT}^{X,\mathrm{Hilb}}_{n,h,d}.
$$
Consequently the physical D-brane partition function is expected to be
$$
Z^X_{\mathrm{phys}}(T,Q,P)
\ \stackrel{\mathrm{phys}}{=}\ 
Z^X_{\mathrm{DT}}(T,Q,P).
$$
Combining this conjectural dictionary with
Theorem~\ref{thm:hilb-dt-scalar-normalization-status} gives the physical
scalar square
$$
Z^X_{\mathrm{phys}}
\ \stackrel{\mathrm{phys}}{=}\ 
\frac{C_{\mathrm{phys}}}{\Delta_5^2}.
$$
The protected-index dictionary is not used in
Theorem~\ref{thm:op-primitive-square-status} or
Theorem~\ref{thm:scalar-square-consequence-status}.
\end{conjecture}

\begin{center}
\small
\begin{tabular}{p{0.20\textwidth}p{0.42\textwidth}p{0.28\textwidth}}
\hline
Object & Equation & Status \\
\hline
OP primitive reduced curve-counting &
\(Z^X_{\mathrm{OP}}=-(\chi_{10}^{\mathrm{OP}})^{-1}
=C_{\mathrm{OP}}\Delta_5^{-2}\) &
Proved by Oberdieck--Pixton for primitive K3 classes. \\
Hilbert-scheme DT normalization &
\(Z^X_{\mathrm{DT}}
=\sum \mathbf{DT}^{X,\mathrm{Hilb}}_{n,h,d}
T^{d-1}Q^{h-1}(-P)^n\) &
Definition.  Equality with \(-4096\Delta_5^{-2}\) follows from
Oberdieck reduced DT/PT and Oberdieck--Pixton. \\
Protected D-brane index &
\(\Omega_X(h,d,n)\stackrel{\mathrm{phys}}{=}
\mathbf{DT}^{X,\mathrm{Hilb}}_{n,h,d}\) &
Conjectural physics dictionary. \\
Scalar-square branch &
\(Z^X_\square=C_\square\Delta_5^{-2}\) &
Unconditional for \(Z^X_\square=Z^X_{\mathrm{OP}}\) and
\(Z^X_\square=Z^X_{\mathrm{DT}}\) in the reduced quotient
normalization. \\
\hline
\end{tabular}
\end{center}
